\documentclass[10pt,a4paper]{amsart}
\usepackage{amsmath,amssymb,amsthm,graphicx}
% Retained minimal preamble for the Stacks Project chapter more-algebra.tex (More on Algebra).
% The chapter file pinned-original-source.tex carries only \input{preamble} on its first line, so no
% preamble text lives inside the chapter itself.  The Stacks Project uses one single shared preamble for
% every chapter; the snapshot stages it identically as sources/stacks-chow/preamble.tex and
% sources/stacks-etale-cohomology/preamble.tex.  The block below is copied mechanically, line for line,
% from sources/stacks-chow/preamble.tex lines 14-18 (the Xy-pic setup) and 154-217 (the native theorem
% declarations together with the mathematical macros).  Nothing was added, removed or redefined; the
% wrapper class and the amsmath/amssymb/amsthm/graphicx base packages are supplied by the assembler.
% Chapter: more-algebra.tex -- More on Algebra.
% For commutative diagrams we use Xy-pic
\usepackage[all]{xy}

% We use 2cell for 2-commutative diagrams.
\xyoption{2cell}

% Theorem environments.
%
\theoremstyle{plain}
\newtheorem{theorem}[subsection]{Theorem}
\newtheorem{proposition}[subsection]{Proposition}
\newtheorem{lemma}[subsection]{Lemma}

\theoremstyle{definition}
\newtheorem{definition}[subsection]{Definition}
\newtheorem{example}[subsection]{Example}
\newtheorem{exercise}[subsection]{Exercise}
\newtheorem{situation}[subsection]{Situation}

\theoremstyle{remark}
\newtheorem{remark}[subsection]{Remark}
\newtheorem{remarks}[subsection]{Remarks}

\numberwithin{equation}{subsection}

% Macros
%
\def\lim{\mathop{\mathrm{lim}}\nolimits}
\def\colim{\mathop{\mathrm{colim}}\nolimits}
\def\Spec{\mathop{\mathrm{Spec}}}
\def\Hom{\mathop{\mathrm{Hom}}\nolimits}
\def\Ext{\mathop{\mathrm{Ext}}\nolimits}
\def\SheafHom{\mathop{\mathcal{H}\!\mathit{om}}\nolimits}
\def\SheafExt{\mathop{\mathcal{E}\!\mathit{xt}}\nolimits}
\def\Sch{\mathit{Sch}}
\def\Mor{\mathop{\mathrm{Mor}}\nolimits}
\def\Ob{\mathop{\mathrm{Ob}}\nolimits}
\def\Sh{\mathop{\mathit{Sh}}\nolimits}
\def\NL{\mathop{N\!L}\nolimits}
\def\CH{\mathop{\mathrm{CH}}\nolimits}
\def\proetale{{pro\text{-}\acute{e}tale}}
\def\etale{{\acute{e}tale}}
\def\QCoh{\mathit{QCoh}}
\def\Ker{\mathop{\mathrm{Ker}}}
\def\Im{\mathop{\mathrm{Im}}}
\def\Coker{\mathop{\mathrm{Coker}}}
\def\Coim{\mathop{\mathrm{Coim}}}

% Boxtimes
%
\DeclareMathSymbol{\boxtimes}{\mathbin}{AMSa}{"02}

%
% Macros for moduli stacks/spaces
%
\def\QCohstack{\mathcal{QC}\!\mathit{oh}}
\def\Cohstack{\mathcal{C}\!\mathit{oh}}
\def\Spacesstack{\mathcal{S}\!\mathit{paces}}
\def\Quotfunctor{\mathrm{Quot}}
\def\Hilbfunctor{\mathrm{Hilb}}
\def\Curvesstack{\mathcal{C}\!\mathit{urves}}
\def\Polarizedstack{\mathcal{P}\!\mathit{olarized}}
\def\Complexesstack{\mathcal{C}\!\mathit{omplexes}}
% \Pic is the operator that assigns to X its picard group, usage \Pic(X)
% \Picardstack_{X/B} denotes the Picard stack of X over B
% \Picardfunctor_{X/B} denotes the Picard functor of X over B
\def\Pic{\mathop{\mathrm{Pic}}\nolimits}
\def\Picardstack{\mathcal{P}\!\mathit{ic}}
\def\Picardfunctor{\mathrm{Pic}}
\def\Deformationcategory{\mathcal{D}\!\mathit{ef}}
\title{A module flat over A is flat over B when B tensor over A B maps flatly onto B}
\author{The Stacks Project Authors}
\date{stacks-project:more-algebra.tex:local-snapshot-20261012; GNU FDL 1.2 or later}
\begin{document}
\maketitle
\setcounter{subsection}{1}
\noindent\textit{Editorial note.} All delivered mathematical text is the author's own, copied from the pinned Stacks Project chapter file more-algebra.tex; no statement or proof step was rewritten, reordered or removed. The counted statement is the labeled lemma environment at source lines 30201-30206 (canonical tag 092C) and the counted proof is the proof environment at source lines 30208-30221, which immediately follows it and is its own proof.\par\medskip
\noindent\textit{Editorial note.} The lead-in sentence at source lines 30198-30199, recording that an etale ring map is weakly etale and announcing that a simple but key property follows, lies before the counted statement and is outside the counted spans.\par\medskip
\noindent\textit{Editorial note.} The chapter file pinned-original-source.tex carries only an input of the shared preamble on its first line, so no preamble text lives inside the chapter. The retained file sources/stacks-more-algebra/retained-preamble.tex is a mechanical, line-for-line copy of the project's single shared preamble as staged in the snapshot (Xy-pic setup plus the native theorem declarations and mathematical macros). Nothing was added, removed or redefined, and the excerpt serves only as preamble, never as a span coordinate.\par\medskip
\noindent\textit{Editorial note.} The source URL is the source-level Stacks Project URL, as required when more than one card is produced from one source; the per-claim tag is carried by claim\_key and by this note. The counted statement and proof refer to no labelled item of the chapter and cite no bibliography entry, so this card contributes no external reference stub and needs no bibliography.\par\medskip
\noindent\textit{Allowed standard prerequisites (editorial).}\begin{itemize}
\item A module N is flat over A if and only if tensoring with N over A is an exact functor on the category of A-modules. used at proof lines 30209-30214, where exactness of the functor sending N' to N tensor\_A N' is deduced from flatness of N over A. Produces an exact functor out of the category of B-modules as the first half of the composite that has to be shown exact.
\item Extension of scalars along a flat ring map is exact, so tensoring with B over B tensor\_A B preserves exactness. used at proof lines 30215-30218, where flatness of B tensor\_A B to B is applied. Keeps exactness when the functor is composed with the change of scalars back to B.
\item A composite of the two functors is naturally isomorphic to tensoring with N over B, so N is flat over B. used at proof lines 30217-30220, where (N tensor\_A N') tensor over B tensor\_A B B is identified with N tensor\_B N'. Identifies the exact functor constructed with the tensor functor of N over B, which is the conclusion.
\end{itemize}\medskip
\begin{lemma}
\label{lemma-key}
Let $A \to B$ be a ring map such that $B \otimes_A B \to B$ is flat.
Let $N$ be a $B$-module. If $N$ is flat as an $A$-module, then
$N$ is flat as a $B$-module.
\end{lemma}

\begin{proof}
Assume $N$ is a flat as an $A$-module.
Then the functor
$$
\text{Mod}_B \longrightarrow \text{Mod}_{B \otimes_A B},\quad
N' \mapsto N \otimes_A N'
$$
is exact. As $B \otimes_A B \to B$ is flat we conclude that the functor
$$
\text{Mod}_B \longrightarrow \text{Mod}_B,\quad
N' \mapsto (N \otimes_A N') \otimes_{B \otimes_A B} B = N \otimes_B N'
$$
is exact, hence $N$ is flat over $B$.
\end{proof}

\end{document}
